\chapter{Traps and system calls}
\label{CH:TRAP}

There are three kinds of event which cause the CPU to set
aside ordinary execution of instructions and force a
transfer of control to special kernel code that handles the event. One
situation is a system call, when a user program 
executes the {\tt ecall} instruction to ask the kernel to do 
something for it. Another situation is an \indextext{exception}:
an instruction (user or kernel) does something illegal, such as 
load from an invalid virtual address. The third situation is a
device \indextext{interrupt}, when a device signals that it needs
attention, for example when the disk hardware finishes a read or write
request.

This book uses \indextext{trap} as a generic term for these
situations. Typically whatever code was executing at the time of the
trap will later need to resume, and shouldn't need to be aware that
anything special happened. That is, we often want traps to be
transparent; this is particularly important for device interrupts, which the
interrupted code typically doesn't expect. 
A trap forces a transfer of control into the kernel; the kernel saves
registers and other state so that execution can be resumed; the kernel
executes appropriate handler code (e.g., a system call implementation
or device driver); the kernel restores the saved state and returns
from the trap; and the original code resumes where it left off.

Xv6 handles all traps in the kernel; traps are not delivered to user
code. Handling traps in the kernel is natural for system calls. It
makes sense for interrupts since isolation demands that only the
kernel be allowed to use devices,
and because the kernel is able to
share devices among multiple processes.
It also makes sense for exceptions since the kernel may be able to
handle the exception from user space (for an example see
Chapter~\ref{CH:PGFAULTS}) or respond by killing the offending
program.

Xv6 trap handling proceeds in four stages: hardware actions taken by
the RISC-V CPU, some assembly instructions that prepare the way for
kernel C code, a C function that decides what to do with the trap,
and the system call or device-driver service routine. While
commonality among the three trap types suggests that a kernel could
handle all traps with a single code path, it turns out to be
convenient to have separate code for
two distinct cases: traps from user space, and traps from kernel space.
Kernel code (assembler or C) that
processes a trap is often called a \indextext{handler};
the first handler instructions are usually written in assembler
(rather than C) and are sometimes called a \indextext{vector}.

Before proceeding, please read \lstinline{kernel/trampoline.S} \fileref{kernel/trampoline.S},
and {\tt usertrap()} and {\tt prepare\_return()} in \lstinline{kernel/trap.c} \fileref{kernel/trap.c}.

\section{RISC-V trap machinery}

Each RISC-V CPU has a set of hardware control registers that the kernel writes to
tell the CPU how to handle traps, and that the kernel can read
to find out about a trap that has occurred. The RISC-V documents
contain the full story~\cite{riscv:priv}. {\tt riscv.h}
\fileref{kernel/riscv.h} contains definitions that xv6 uses. Here's
an outline of the most important registers:

\begin{itemize}

\item \indexcode{stvec}: The kernel writes the address of its trap handler
  code here; the RISC-V jumps to the address in {\tt stvec} to handle a trap.

\item \indexcode{sepc}: When a trap occurs, RISC-V saves the program counter
  here (since the {\tt pc} is then overwritten with the
  value in {\tt stvec}). The
  {\tt sret} (return from trap) instruction copies {\tt sepc} to the
  {\tt pc}. The kernel can write {\tt sepc} to control where {\tt
    sret} goes.

\item \indexcode{scause}: RISC-V puts a number here that describes
the reason for the trap.

\item \indexcode{sscratch}: The kernel trap handler code uses {\tt sscratch}
  to help it avoid overwriting user registers before saving them.

\item \indexcode{sstatus}: The SIE bit in {\tt sstatus}
  controls whether device interrupts
  are enabled. If the kernel clears SIE, the RISC-V will defer
  device interrupts until the kernel sets SIE. The SPP bit
  indicates whether a trap came from user mode or supervisor
  mode, and controls to what mode {\tt sret} returns.

\end{itemize}

The above registers can only be accessed in supervisor mode (i.e., by
the kernel); the CPU prevents user code from reading or writing them.

Each CPU on a multi-core chip has its own set of these registers,
and more than one CPU may be handling a trap at any given time.

When it forces a trap, the RISC-V hardware does the following:

\begin{enumerate}

\item If the trap is a device interrupt, and the {\tt sstatus} SIE bit
  is clear, don't do any of the following.

\item Disable interrupts by clearing the SIE bit in {\tt sstatus}.

\item Copy the {\tt pc} to {\tt sepc}.

\item Save the current mode (user or supervisor) in the SPP bit in {\tt sstatus}.

\item Set {\tt scause} to a number indicating the trap's cause.

\item Set the mode to supervisor.

\item Copy {\tt stvec} to the {\tt pc}.

\item Start executing at the new {\tt pc}.

\end{enumerate}

The CPU doesn't switch to the kernel page table, doesn't
switch to a stack in the kernel, and doesn't save any registers other
than the {\tt pc}. Kernel software must perform these tasks.
One reason that the CPU does minimal work during a trap is to provide
flexibility to software; for example, some operating systems 
omit a page table switch in some situations to increase
trap performance.

It's worth thinking about whether any of the steps listed above could
be omitted, perhaps in search of faster traps. Though there are
situations in which a simpler sequence can work, many of the steps
would be dangerous to omit in general. For example, suppose that the
CPU didn't switch program counters. Then a trap from user space could
switch to supervisor mode while still running user instructions. Those
user instructions could break user/kernel isolation, for example
by modifying the {\tt satp} register to point to a page table that
allowed accessing all of physical memory. It is thus important that
the CPU switch to a kernel-specified instruction address, namely {\tt
  stvec}.

\section{Traps from user space}
\begin{figure}[t]
  \begin{center}
    \chapter{Traps and system calls}
\label{CH:TRAP}

There are three kinds of event which cause the CPU to set
aside ordinary execution of instructions and force a
transfer of control to special kernel code that handles the event. One
situation is a system call, when a user program 
executes the {\tt ecall} instruction to ask the kernel to do 
something for it. Another situation is an \indextext{exception}:
an instruction (user or kernel) does something illegal, such as 
load from an invalid virtual address. The third situation is a
device \indextext{interrupt}, when a device signals that it needs
attention, for example when the disk hardware finishes a read or write
request.

This book uses \indextext{trap} as a generic term for these
situations. Typically whatever code was executing at the time of the
trap will later need to resume, and shouldn't need to be aware that
anything special happened. That is, we often want traps to be
transparent; this is particularly important for device interrupts, which the
interrupted code typically doesn't expect. 
A trap forces a transfer of control into the kernel; the kernel saves
registers and other state so that execution can be resumed; the kernel
executes appropriate handler code (e.g., a system call implementation
or device driver); the kernel restores the saved state and returns
from the trap; and the original code resumes where it left off.

Xv6 handles all traps in the kernel; traps are not delivered to user
code. Handling traps in the kernel is natural for system calls. It
makes sense for interrupts since isolation demands that only the
kernel be allowed to use devices,
and because the kernel is able to
share devices among multiple processes.
It also makes sense for exceptions since the kernel may be able to
handle the exception from user space (for an example see
Chapter~\ref{CH:PGFAULTS}) or respond by killing the offending
program.

Xv6 trap handling proceeds in four stages: hardware actions taken by
the RISC-V CPU, some assembly instructions that prepare the way for
kernel C code, a C function that decides what to do with the trap,
and the system call or device-driver service routine. While
commonality among the three trap types suggests that a kernel could
handle all traps with a single code path, it turns out to be
convenient to have separate code for
two distinct cases: traps from user space, and traps from kernel space.
Kernel code (assembler or C) that
processes a trap is often called a \indextext{handler};
the first handler instructions are usually written in assembler
(rather than C) and are sometimes called a \indextext{vector}.

Before proceeding, please read \lstinline{kernel/trampoline.S} \fileref{kernel/trampoline.S},
and {\tt usertrap()} and {\tt prepare\_return()} in \lstinline{kernel/trap.c} \fileref{kernel/trap.c}.

\section{RISC-V trap machinery}

Each RISC-V CPU has a set of hardware control registers that the kernel writes to
tell the CPU how to handle traps, and that the kernel can read
to find out about a trap that has occurred. The RISC-V documents
contain the full story~\cite{riscv:priv}. {\tt riscv.h}
\fileref{kernel/riscv.h} contains definitions that xv6 uses. Here's
an outline of the most important registers:

\begin{itemize}

\item \indexcode{stvec}: The kernel writes the address of its trap handler
  code here; the RISC-V jumps to the address in {\tt stvec} to handle a trap.

\item \indexcode{sepc}: When a trap occurs, RISC-V saves the program counter
  here (since the {\tt pc} is then overwritten with the
  value in {\tt stvec}). The
  {\tt sret} (return from trap) instruction copies {\tt sepc} to the
  {\tt pc}. The kernel can write {\tt sepc} to control where {\tt
    sret} goes.

\item \indexcode{scause}: RISC-V puts a number here that describes
the reason for the trap.

\item \indexcode{sscratch}: The kernel trap handler code uses {\tt sscratch}
  to help it avoid overwriting user registers before saving them.

\item \indexcode{sstatus}: The SIE bit in {\tt sstatus}
  controls whether device interrupts
  are enabled. If the kernel clears SIE, the RISC-V will defer
  device interrupts until the kernel sets SIE. The SPP bit
  indicates whether a trap came from user mode or supervisor
  mode, and controls to what mode {\tt sret} returns.

\end{itemize}

The above registers can only be accessed in supervisor mode (i.e., by
the kernel); the CPU prevents user code from reading or writing them.

Each CPU on a multi-core chip has its own set of these registers,
and more than one CPU may be handling a trap at any given time.

When it forces a trap, the RISC-V hardware does the following:

\begin{enumerate}

\item If the trap is a device interrupt, and the {\tt sstatus} SIE bit
  is clear, don't do any of the following.

\item Disable interrupts by clearing the SIE bit in {\tt sstatus}.

\item Copy the {\tt pc} to {\tt sepc}.

\item Save the current mode (user or supervisor) in the SPP bit in {\tt sstatus}.

\item Set {\tt scause} to a number indicating the trap's cause.

\item Set the mode to supervisor.

\item Copy {\tt stvec} to the {\tt pc}.

\item Start executing at the new {\tt pc}.

\end{enumerate}

The CPU doesn't switch to the kernel page table, doesn't
switch to a stack in the kernel, and doesn't save any registers other
than the {\tt pc}. Kernel software must perform these tasks.
One reason that the CPU does minimal work during a trap is to provide
flexibility to software; for example, some operating systems 
omit a page table switch in some situations to increase
trap performance.

It's worth thinking about whether any of the steps listed above could
be omitted, perhaps in search of faster traps. Though there are
situations in which a simpler sequence can work, many of the steps
would be dangerous to omit in general. For example, suppose that the
CPU didn't switch program counters. Then a trap from user space could
switch to supervisor mode while still running user instructions. Those
user instructions could break user/kernel isolation, for example
by modifying the {\tt satp} register to point to a page table that
allowed accessing all of physical memory. It is thus important that
the CPU switch to a kernel-specified instruction address, namely {\tt
  stvec}.

\section{Traps from user space}
\begin{figure}[t]
  \begin{center}
    \chapter{Traps and system calls}
\label{CH:TRAP}

There are three kinds of event which cause the CPU to set
aside ordinary execution of instructions and force a
transfer of control to special kernel code that handles the event. One
situation is a system call, when a user program 
executes the {\tt ecall} instruction to ask the kernel to do 
something for it. Another situation is an \indextext{exception}:
an instruction (user or kernel) does something illegal, such as 
load from an invalid virtual address. The third situation is a
device \indextext{interrupt}, when a device signals that it needs
attention, for example when the disk hardware finishes a read or write
request.

This book uses \indextext{trap} as a generic term for these
situations. Typically whatever code was executing at the time of the
trap will later need to resume, and shouldn't need to be aware that
anything special happened. That is, we often want traps to be
transparent; this is particularly important for device interrupts, which the
interrupted code typically doesn't expect. 
A trap forces a transfer of control into the kernel; the kernel saves
registers and other state so that execution can be resumed; the kernel
executes appropriate handler code (e.g., a system call implementation
or device driver); the kernel restores the saved state and returns
from the trap; and the original code resumes where it left off.

Xv6 handles all traps in the kernel; traps are not delivered to user
code. Handling traps in the kernel is natural for system calls. It
makes sense for interrupts since isolation demands that only the
kernel be allowed to use devices,
and because the kernel is able to
share devices among multiple processes.
It also makes sense for exceptions since the kernel may be able to
handle the exception from user space (for an example see
Chapter~\ref{CH:PGFAULTS}) or respond by killing the offending
program.

Xv6 trap handling proceeds in four stages: hardware actions taken by
the RISC-V CPU, some assembly instructions that prepare the way for
kernel C code, a C function that decides what to do with the trap,
and the system call or device-driver service routine. While
commonality among the three trap types suggests that a kernel could
handle all traps with a single code path, it turns out to be
convenient to have separate code for
two distinct cases: traps from user space, and traps from kernel space.
Kernel code (assembler or C) that
processes a trap is often called a \indextext{handler};
the first handler instructions are usually written in assembler
(rather than C) and are sometimes called a \indextext{vector}.

Before proceeding, please read \lstinline{kernel/trampoline.S} \fileref{kernel/trampoline.S},
and {\tt usertrap()} and {\tt prepare\_return()} in \lstinline{kernel/trap.c} \fileref{kernel/trap.c}.

\section{RISC-V trap machinery}

Each RISC-V CPU has a set of hardware control registers that the kernel writes to
tell the CPU how to handle traps, and that the kernel can read
to find out about a trap that has occurred. The RISC-V documents
contain the full story~\cite{riscv:priv}. {\tt riscv.h}
\fileref{kernel/riscv.h} contains definitions that xv6 uses. Here's
an outline of the most important registers:

\begin{itemize}

\item \indexcode{stvec}: The kernel writes the address of its trap handler
  code here; the RISC-V jumps to the address in {\tt stvec} to handle a trap.

\item \indexcode{sepc}: When a trap occurs, RISC-V saves the program counter
  here (since the {\tt pc} is then overwritten with the
  value in {\tt stvec}). The
  {\tt sret} (return from trap) instruction copies {\tt sepc} to the
  {\tt pc}. The kernel can write {\tt sepc} to control where {\tt
    sret} goes.

\item \indexcode{scause}: RISC-V puts a number here that describes
the reason for the trap.

\item \indexcode{sscratch}: The kernel trap handler code uses {\tt sscratch}
  to help it avoid overwriting user registers before saving them.

\item \indexcode{sstatus}: The SIE bit in {\tt sstatus}
  controls whether device interrupts
  are enabled. If the kernel clears SIE, the RISC-V will defer
  device interrupts until the kernel sets SIE. The SPP bit
  indicates whether a trap came from user mode or supervisor
  mode, and controls to what mode {\tt sret} returns.

\end{itemize}

The above registers can only be accessed in supervisor mode (i.e., by
the kernel); the CPU prevents user code from reading or writing them.

Each CPU on a multi-core chip has its own set of these registers,
and more than one CPU may be handling a trap at any given time.

When it forces a trap, the RISC-V hardware does the following:

\begin{enumerate}

\item If the trap is a device interrupt, and the {\tt sstatus} SIE bit
  is clear, don't do any of the following.

\item Disable interrupts by clearing the SIE bit in {\tt sstatus}.

\item Copy the {\tt pc} to {\tt sepc}.

\item Save the current mode (user or supervisor) in the SPP bit in {\tt sstatus}.

\item Set {\tt scause} to a number indicating the trap's cause.

\item Set the mode to supervisor.

\item Copy {\tt stvec} to the {\tt pc}.

\item Start executing at the new {\tt pc}.

\end{enumerate}

The CPU doesn't switch to the kernel page table, doesn't
switch to a stack in the kernel, and doesn't save any registers other
than the {\tt pc}. Kernel software must perform these tasks.
One reason that the CPU does minimal work during a trap is to provide
flexibility to software; for example, some operating systems 
omit a page table switch in some situations to increase
trap performance.

It's worth thinking about whether any of the steps listed above could
be omitted, perhaps in search of faster traps. Though there are
situations in which a simpler sequence can work, many of the steps
would be dangerous to omit in general. For example, suppose that the
CPU didn't switch program counters. Then a trap from user space could
switch to supervisor mode while still running user instructions. Those
user instructions could break user/kernel isolation, for example
by modifying the {\tt satp} register to point to a page table that
allowed accessing all of physical memory. It is thus important that
the CPU switch to a kernel-specified instruction address, namely {\tt
  stvec}.

\section{Traps from user space}
\begin{figure}[t]
  \begin{center}
    \chapter{Traps and system calls}
\label{CH:TRAP}

There are three kinds of event which cause the CPU to set
aside ordinary execution of instructions and force a
transfer of control to special kernel code that handles the event. One
situation is a system call, when a user program 
executes the {\tt ecall} instruction to ask the kernel to do 
something for it. Another situation is an \indextext{exception}:
an instruction (user or kernel) does something illegal, such as 
load from an invalid virtual address. The third situation is a
device \indextext{interrupt}, when a device signals that it needs
attention, for example when the disk hardware finishes a read or write
request.

This book uses \indextext{trap} as a generic term for these
situations. Typically whatever code was executing at the time of the
trap will later need to resume, and shouldn't need to be aware that
anything special happened. That is, we often want traps to be
transparent; this is particularly important for device interrupts, which the
interrupted code typically doesn't expect. 
A trap forces a transfer of control into the kernel; the kernel saves
registers and other state so that execution can be resumed; the kernel
executes appropriate handler code (e.g., a system call implementation
or device driver); the kernel restores the saved state and returns
from the trap; and the original code resumes where it left off.

Xv6 handles all traps in the kernel; traps are not delivered to user
code. Handling traps in the kernel is natural for system calls. It
makes sense for interrupts since isolation demands that only the
kernel be allowed to use devices,
and because the kernel is able to
share devices among multiple processes.
It also makes sense for exceptions since the kernel may be able to
handle the exception from user space (for an example see
Chapter~\ref{CH:PGFAULTS}) or respond by killing the offending
program.

Xv6 trap handling proceeds in four stages: hardware actions taken by
the RISC-V CPU, some assembly instructions that prepare the way for
kernel C code, a C function that decides what to do with the trap,
and the system call or device-driver service routine. While
commonality among the three trap types suggests that a kernel could
handle all traps with a single code path, it turns out to be
convenient to have separate code for
two distinct cases: traps from user space, and traps from kernel space.
Kernel code (assembler or C) that
processes a trap is often called a \indextext{handler};
the first handler instructions are usually written in assembler
(rather than C) and are sometimes called a \indextext{vector}.

Before proceeding, please read \lstinline{kernel/trampoline.S} \fileref{kernel/trampoline.S},
and {\tt usertrap()} and {\tt prepare\_return()} in \lstinline{kernel/trap.c} \fileref{kernel/trap.c}.

\section{RISC-V trap machinery}

Each RISC-V CPU has a set of hardware control registers that the kernel writes to
tell the CPU how to handle traps, and that the kernel can read
to find out about a trap that has occurred. The RISC-V documents
contain the full story~\cite{riscv:priv}. {\tt riscv.h}
\fileref{kernel/riscv.h} contains definitions that xv6 uses. Here's
an outline of the most important registers:

\begin{itemize}

\item \indexcode{stvec}: The kernel writes the address of its trap handler
  code here; the RISC-V jumps to the address in {\tt stvec} to handle a trap.

\item \indexcode{sepc}: When a trap occurs, RISC-V saves the program counter
  here (since the {\tt pc} is then overwritten with the
  value in {\tt stvec}). The
  {\tt sret} (return from trap) instruction copies {\tt sepc} to the
  {\tt pc}. The kernel can write {\tt sepc} to control where {\tt
    sret} goes.

\item \indexcode{scause}: RISC-V puts a number here that describes
the reason for the trap.

\item \indexcode{sscratch}: The kernel trap handler code uses {\tt sscratch}
  to help it avoid overwriting user registers before saving them.

\item \indexcode{sstatus}: The SIE bit in {\tt sstatus}
  controls whether device interrupts
  are enabled. If the kernel clears SIE, the RISC-V will defer
  device interrupts until the kernel sets SIE. The SPP bit
  indicates whether a trap came from user mode or supervisor
  mode, and controls to what mode {\tt sret} returns.

\end{itemize}

The above registers can only be accessed in supervisor mode (i.e., by
the kernel); the CPU prevents user code from reading or writing them.

Each CPU on a multi-core chip has its own set of these registers,
and more than one CPU may be handling a trap at any given time.

When it forces a trap, the RISC-V hardware does the following:

\begin{enumerate}

\item If the trap is a device interrupt, and the {\tt sstatus} SIE bit
  is clear, don't do any of the following.

\item Disable interrupts by clearing the SIE bit in {\tt sstatus}.

\item Copy the {\tt pc} to {\tt sepc}.

\item Save the current mode (user or supervisor) in the SPP bit in {\tt sstatus}.

\item Set {\tt scause} to a number indicating the trap's cause.

\item Set the mode to supervisor.

\item Copy {\tt stvec} to the {\tt pc}.

\item Start executing at the new {\tt pc}.

\end{enumerate}

The CPU doesn't switch to the kernel page table, doesn't
switch to a stack in the kernel, and doesn't save any registers other
than the {\tt pc}. Kernel software must perform these tasks.
One reason that the CPU does minimal work during a trap is to provide
flexibility to software; for example, some operating systems 
omit a page table switch in some situations to increase
trap performance.

It's worth thinking about whether any of the steps listed above could
be omitted, perhaps in search of faster traps. Though there are
situations in which a simpler sequence can work, many of the steps
would be dangerous to omit in general. For example, suppose that the
CPU didn't switch program counters. Then a trap from user space could
switch to supervisor mode while still running user instructions. Those
user instructions could break user/kernel isolation, for example
by modifying the {\tt satp} register to point to a page table that
allowed accessing all of physical memory. It is thus important that
the CPU switch to a kernel-specified instruction address, namely {\tt
  stvec}.

\section{Traps from user space}
\begin{figure}[t]
  \begin{center}
    \input{fig/trap.tex}
  \end{center}
\caption{Outline of how a trap from user code is handled.}
\label{fig:usertrap}
\end{figure}

Xv6 handles traps differently depending on whether
the trap occurs while
executing in the kernel or in user code. Here is the
story for traps from user code; Section~\ref{s:ktraps}
describes traps from kernel code.

A trap may occur while executing in user space if the
user program makes a
system call ({\tt ecall} instruction), or does something
illegal, or if a device interrupts.
As shown in Figure~\ref{fig:usertrap}, the high-level path of a trap from user space is
{\tt uservec}
\lineref{kernel/trampoline.S:/^uservec/},
then {\tt usertrap}
\lineref{kernel/trap.c:/^usertrap/};
and when the kernel is ready to return,
{\tt usertrap}
returns to
{\tt userret}
\lineref{kernel/trampoline.S:/^userret/}
which executes {\tt sret}
to user space.

% talk about why RISC-V doesn't switch page tables

A major constraint on the design of xv6's trap handling is the fact
that the RISC-V hardware does not switch page tables when it forces a
trap. This means that the trap handler
address in {\tt stvec} must have a valid
mapping in the user page table, since that's the page table in force
when the trap handling code starts executing. Furthermore, xv6's trap
handling code needs to switch to the kernel page table; in order to be
able to continue executing after that switch, the kernel page table
must also have a mapping for the handler pointed to by {\tt stvec}.

Xv6 satisfies these requirements using a \indextext{trampoline} page.
This page contains {\tt uservec}, the xv6 trap handling code
that {\tt stvec} points to. The trampoline page is mapped in
every process's page table at virtual address {\tt 0x3ffffff000}
(called \indexcode{TRAMPOLINE}),
which is the last page in the virtual address space so that it will
be above memory that programs use for themselves.
The trampoline page is mapped at the same virtual address
in the kernel page table. See Figure~\ref{fig:as} and
Figure~\ref{fig:xv6_layout}. Because the trampoline page is mapped in
the user page table, traps can start
executing there in supervisor mode. Because the trampoline page is
mapped at the same address in the kernel address space, the trap handler
can continue to execute after it switches to the kernel page
table.

The code for the {\tt uservec} trap handler is in {\tt trampoline.S}
\lineref{kernel/trampoline.S:/^uservec/}.
When {\tt uservec} starts, all 32 registers contain values owned by
the interrupted user code. These 32 values need to be saved somewhere
in memory, so that later on
the kernel can restore them before returning to user space.
Storing to memory requires use of a register
to hold the store's destination address,
but at this point there are no general-purpose registers available!
Luckily RISC-V provides a helping hand in the
form of the {\tt sscratch} register. The {\tt csrw} instruction at
the start of {\tt uservec} saves {\tt a0} in {\tt
  sscratch}. Now 
{\tt uservec} has
one register ({\tt a0}) to play with.

{\tt uservec}'s next task is to save the 32 user registers.
The kernel allocates, for each process, a page of memory for a
{\tt trapframe} structure that (among other things) has space to
save the 32 user registers
\lineref{kernel/proc.h:/^struct.trapframe/}. Because {\tt satp} still
refers to the user page table, {\tt uservec} needs the trapframe to be
mapped in the user address space. Xv6 maps each process's trapframe
at virtual address {\tt TRAPFRAME} ({\tt 0x3fffffe000}) in that process's user page table;
one page below {\tt TRAMPOLINE}.
Each process's {\tt p->trapframe} 
contains a kernel virtual address for the process's trapframe.

{\tt uservec} sets register {\tt a0} to address {\tt TRAPFRAME}
and saves all the user registers there.
Then it retrieves the user {\tt a0} from {\tt sscratch} and
saves it in the trapframe.

The kernel previously initialized the trapframe to contain some
values useful to {\tt uservec}:
the address of the current process's
kernel stack, the current CPU's hartid, the address of the {\tt usertrap}
function,
and the address of the kernel page table. {\tt uservec}
retrieves these values, switches {\tt satp} to the kernel page table,
and jumps to {\tt usertrap}, a C function.

The job of {\tt usertrap} is to determine
the cause of the trap, process it, and return
\lineref{kernel/trap.c:/^usertrap/}.
It first changes {\tt stvec} so
that a trap while in the kernel will be handled by
{\tt kernelvec} rather than {\tt uservec}.
It saves the {\tt sepc} register (the saved user program counter)
for future use when returning back to user space.
If the trap is a system call, {\tt usertrap} calls {\tt syscall} to
handle it;
if a device interrupt, {\tt devintr};
if a page fault,  {\tt vmfault};
otherwise it's an exception (e.g., use of an invalid address),
and the kernel kills the faulting process.
The system call path adds four to the saved user program counter
because RISC-V, in the case of a system call,
leaves the program pointer pointing to the {\tt ecall} instruction
but user code needs to resume executing at the subsequent instruction.
{\tt usertrap} checks if the process has been
killed or should yield the CPU (if this trap is a timer interrupt).

The first step in returning to user space is the call to {\tt prepare\_return}
{\lineref{kernel/trap.c:/^prepare/}}.
This function sets up the RISC-V control registers to prepare for a
future trap from user space: setting {\tt stvec}
to {\tt uservec} and preparing the trapframe fields that
{\tt uservec} relies on.
{\tt prepare\_return} sets {\tt sepc} to the previously
saved user program counter.
Finally, {\tt usertrap}
returns back to {\tt userret} in the trampoline page
\lineref{kernel/trampoline.S:/^userret/},
passing back a pointer to the user page table in {\tt a0}.

{\tt userret} switches {\tt satp} to the process's user page table.
Recall that the user page table maps both the trampoline page
and {\tt TRAPFRAME}, but nothing else from the kernel.
The trampoline page mapping at the same
virtual address in user and kernel page tables allows
{\tt userret} to keep executing after changing {\tt satp}.
From this point on, the only data {\tt userret} can use is
the register contents and the content of the trapframe.
{\tt userret} loads the {\tt TRAPFRAME} address into {\tt a0},
restores saved user registers from the trapframe via {\tt a0},
restores the saved user {\tt a0},
and executes {\tt sret} to return to user space.

{\tt uservec} and {\tt userret} are written in assembly
language because it is difficult to write C code to
save or  restore all the registers or survive switching
page tables.

\section{Code: Calling system calls}

User programs call library functions in order to make system calls.
For example, the shell displays a prompt with this function call
(in {\tt user/sh.c}):

\begin{verbatim}
  write(2, "$ ", 2);
\end{verbatim}

Here's the library function, in {\tt user/usys.S}:

\begin{verbatim}
write:
 li a7, SYS_write
 ecall
 ret
\end{verbatim}

The code that the C compiler generates for the function call
loads the three arguments into registers {\tt a0},
{\tt a1}, and {\tt a2}. Then the {\tt write()} function
loads the system call number, {\tt SYS\_write} (16), into {\tt a7}.
The kernel will look at those registers to find out what
system call is intended, and what the arguments are.
The \lstinline{ecall} instruction traps from user space
into the kernel and causes
{\tt uservec},
{\tt usertrap}, and then {\tt syscall} to execute.

At this point, please read \lstinline{kernel/syscall.c} \fileref{kernel/syscall.c},
{\tt sys\_write()} in \lstinline{kernel/sysfile.c} \fileref{kernel/sysfile.c},
and {\tt copyout()}, {\tt copyin()}, and {\tt copyinstr()} in \lstinline{kernel/vm.c} \fileref{kernel/vm.c}.

\indexcode{syscall}
\lineref{kernel/syscall.c:/^syscall/} 
retrieves the system call number from the saved
\texttt{a7} in the trapframe
and uses it to index into {\tt syscalls}
\lineref{kernel/syscall.c:/syscalls/}.
For our example, 
\texttt{a7}
contains 
\indexcode{SYS_write}
\lineref{kernel/syscall.h:/SYS_write/},
resulting in a call to the system call implementation function
\lstinline{sys_write}.

When \lstinline{sys_write} returns,
\lstinline{syscall}
records its return value in
\lstinline{p->trapframe->a0}.
This will cause the original user-space call to 
{\tt write()} to return that value, since the C
calling convention on RISC-V places return values in {\tt a0}.
System calls conventionally return negative numbers to indicate
errors, and zero or positive numbers for success.

\section{Code: System call arguments}

System call arguments start out in the user registers, and
are then moved to the trap frame by the kernel trap code.
The kernel functions
\lstinline{argint},
\lstinline{argaddr},
and
\lstinline{argfd}
retrieve the 
\textit{n} 'th 
system call argument
from the trap frame
as an integer, pointer, or a file descriptor.

Some system calls pass pointers as arguments, and the kernel must use
those pointers to read or write user memory. The {\tt write} system
call, for example, passes the kernel a user-space pointer
to the data to be written.
Such pointers pose
two challenges. First, the user program may be buggy or malicious, and
may pass the kernel an invalid pointer or a pointer intended to trick
the kernel into accessing kernel memory instead of user memory.
Second, the xv6 kernel page table mappings are not the same as the
user page table mappings, so the kernel cannot use ordinary
instructions to load or store from user-supplied addresses.

The kernel implements functions that safely transfer data to and
from user-supplied addresses.
{\tt fetchstr} is an example \lineref{kernel/syscall.c:/^fetchstr/}.
File system calls such as
{\tt exec} use {\tt fetchstr} to retrieve string file-name arguments from user
space.
\lstinline{fetchstr} calls \lstinline{copyinstr}
to do the hard work.

\indexcode{copyinstr}
\lineref{kernel/vm.c:/^copyinstr/} copies up to \lstinline{max} bytes to
\lstinline{dst} from virtual address \lstinline{srcva} in the user page
table \lstinline{pagetable}.
Since \lstinline{pagetable} is {\it not} the current page
table,
\lstinline{copyinstr} uses {\tt walkaddr}
(which calls {\tt walk}) to look up
\lstinline{srcva} in
\lstinline{pagetable}, yielding
physical address \lstinline{pa0}.
The kernel's page table maps all of physical RAM 
at virtual addresses that are equal to the RAM's physical address.
This allows
{\tt copyinstr} to directly copy string bytes from {\tt pa0} to {\tt dst}.
{\tt walkaddr} 
\lineref{kernel/vm.c:/^walkaddr/}
checks that the user-supplied virtual address is part of
the process's user address space, so programs
cannot trick the kernel into reading other memory.
A similar function, {\tt copyout}, copies data from the
kernel to a user-supplied address.

\section{Traps from kernel space}
\label{s:ktraps}

Please read \lstinline{kernel/kernelvec.S} \fileref{kernel/kernelvec.S},
and {\tt kerneltrap()} in \lstinline{kernel/trap.c} \fileref{kernel/trap.c}.

Xv6 handles traps from kernel code in a different way
than traps from user code.
When entering the kernel, {\tt usertrap} points {\tt stvec}
to the assembly code at {\tt kernelvec}
\lineref{kernel/kernelvec.S:/^kernelvec/}.
Since {\tt kernelvec} only executes if
xv6 was already in the kernel, {\tt kernelvec} can rely
on {\tt satp} being set to the kernel page table, and on the
stack pointer referring to a valid kernel stack.
{\tt kernelvec} pushes all 32 registers onto the current stack,
from which it will later restore them
so that the interrupted
kernel code can resume without disturbance.

{\tt kernelvec} saves the registers on the stack of the interrupted
kernel thread, which makes sense because the register values belong to
that thread. This is particularly important if the trap causes a
switch to a different thread -- in that case the trap will actually
return from the stack of the new thread, leaving the interrupted
thread's saved registers safely on its stack.

{\tt kernelvec} jumps to {\tt kerneltrap}
\lineref{kernel/trap.c:/^kerneltrap/} after saving registers.
{\tt kerneltrap} is prepared for just one type of trap:
device interrupts. It calls
{\tt devintr}
\lineref{kernel/trap.c:/^devintr/}
to handle them.
If the trap isn't a device interrupt, it must be an exception,
such as kernel code trying to use an invalid pointer.
This could only be caused by a bug in the kernel code.
The kernel does not have a way to recover in this situation,
so it calls {\tt panic()}, which prints an error message
and then halts.

If {\tt kerneltrap} was called due to a timer interrupt, and a
process's kernel thread is running (as opposed to a scheduler thread),
{\tt kerneltrap} calls {\tt yield} to give other threads a chance to
run. At some point one of those threads will yield, and let our thread
and its {\tt kerneltrap} resume again.
Chapter~\ref{CH:SCHED} explains what happens in {\tt yield}.

When {\tt kerneltrap}'s work is done, it needs to return to whatever
code was interrupted by the trap. Because a {\tt yield} may have
disturbed {\tt sepc} and the previous mode in {\tt sstatus},
{\tt kerneltrap} saves them when it starts. It now restores those
control registers and returns to {\tt kernelvec}
\lineref{kernel/kernelvec.S:/call.kerneltrap$/}.
{\tt kernelvec} pops the saved registers from the stack and
executes {\tt sret}, which copies {\tt sepc} to {\tt pc}
and resumes the interrupted kernel code.

Xv6 sets a CPU's {\tt stvec} to {\tt kernelvec} when that CPU
enters the kernel from user space; you can see this in {\tt usertrap}
\lineref{kernel/trap.c:/DOC: kernelvec/}.
But there's a window of time when the kernel has started executing
but {\tt stvec} is still set to {\tt uservec}, and it's crucial that 
no device interrupt occur during that window.
Luckily the RISC-V always disables interrupts when it starts
to take a trap, and {\tt usertrap} doesn't enable them again until
after it sets {\tt stvec}.

\section{Real world}

The need for trampoline pages could be eliminated if kernel
memory were mapped into every process's user page table (with
\lstinline{PTE_U} clear).
That would
also eliminate the need for a page table switch when trapping from
user space into the kernel. That in turn would allow system call
implementations in the kernel to take advantage of the current
process's user memory being mapped, allowing kernel code to directly
dereference user pointers. Many operating systems have used these ideas to
increase efficiency. Xv6 avoids them in order to reduce the chances of
security bugs in the kernel due to inadvertent use of user pointers,
and to reduce some complexity that would be required to ensure that
user and kernel virtual addresses don't overlap.

\section{Exercises}

\begin{enumerate}

\item Could some or all of the code in {\tt trampoline.S} and
{\tt kernelvec.S} be written in C rather than assembler?

\item Is there a way to eliminate the special {\tt
  TRAPFRAME} page mapping in every user address space? For
  example, could
  {\tt uservec} be modified to simply push the 32 user registers
  onto the kernel stack, or store them in the {\tt proc}
  structure?

\item Could xv6 be modified to eliminate the special {\tt
  TRAMPOLINE} page mapping?

\end{enumerate}

  \end{center}
\caption{Outline of how a trap from user code is handled.}
\label{fig:usertrap}
\end{figure}

Xv6 handles traps differently depending on whether
the trap occurs while
executing in the kernel or in user code. Here is the
story for traps from user code; Section~\ref{s:ktraps}
describes traps from kernel code.

A trap may occur while executing in user space if the
user program makes a
system call ({\tt ecall} instruction), or does something
illegal, or if a device interrupts.
As shown in Figure~\ref{fig:usertrap}, the high-level path of a trap from user space is
{\tt uservec}
\lineref{kernel/trampoline.S:/^uservec/},
then {\tt usertrap}
\lineref{kernel/trap.c:/^usertrap/};
and when the kernel is ready to return,
{\tt usertrap}
returns to
{\tt userret}
\lineref{kernel/trampoline.S:/^userret/}
which executes {\tt sret}
to user space.

% talk about why RISC-V doesn't switch page tables

A major constraint on the design of xv6's trap handling is the fact
that the RISC-V hardware does not switch page tables when it forces a
trap. This means that the trap handler
address in {\tt stvec} must have a valid
mapping in the user page table, since that's the page table in force
when the trap handling code starts executing. Furthermore, xv6's trap
handling code needs to switch to the kernel page table; in order to be
able to continue executing after that switch, the kernel page table
must also have a mapping for the handler pointed to by {\tt stvec}.

Xv6 satisfies these requirements using a \indextext{trampoline} page.
This page contains {\tt uservec}, the xv6 trap handling code
that {\tt stvec} points to. The trampoline page is mapped in
every process's page table at virtual address {\tt 0x3ffffff000}
(called \indexcode{TRAMPOLINE}),
which is the last page in the virtual address space so that it will
be above memory that programs use for themselves.
The trampoline page is mapped at the same virtual address
in the kernel page table. See Figure~\ref{fig:as} and
Figure~\ref{fig:xv6_layout}. Because the trampoline page is mapped in
the user page table, traps can start
executing there in supervisor mode. Because the trampoline page is
mapped at the same address in the kernel address space, the trap handler
can continue to execute after it switches to the kernel page
table.

The code for the {\tt uservec} trap handler is in {\tt trampoline.S}
\lineref{kernel/trampoline.S:/^uservec/}.
When {\tt uservec} starts, all 32 registers contain values owned by
the interrupted user code. These 32 values need to be saved somewhere
in memory, so that later on
the kernel can restore them before returning to user space.
Storing to memory requires use of a register
to hold the store's destination address,
but at this point there are no general-purpose registers available!
Luckily RISC-V provides a helping hand in the
form of the {\tt sscratch} register. The {\tt csrw} instruction at
the start of {\tt uservec} saves {\tt a0} in {\tt
  sscratch}. Now 
{\tt uservec} has
one register ({\tt a0}) to play with.

{\tt uservec}'s next task is to save the 32 user registers.
The kernel allocates, for each process, a page of memory for a
{\tt trapframe} structure that (among other things) has space to
save the 32 user registers
\lineref{kernel/proc.h:/^struct.trapframe/}. Because {\tt satp} still
refers to the user page table, {\tt uservec} needs the trapframe to be
mapped in the user address space. Xv6 maps each process's trapframe
at virtual address {\tt TRAPFRAME} ({\tt 0x3fffffe000}) in that process's user page table;
one page below {\tt TRAMPOLINE}.
Each process's {\tt p->trapframe} 
contains a kernel virtual address for the process's trapframe.

{\tt uservec} sets register {\tt a0} to address {\tt TRAPFRAME}
and saves all the user registers there.
Then it retrieves the user {\tt a0} from {\tt sscratch} and
saves it in the trapframe.

The kernel previously initialized the trapframe to contain some
values useful to {\tt uservec}:
the address of the current process's
kernel stack, the current CPU's hartid, the address of the {\tt usertrap}
function,
and the address of the kernel page table. {\tt uservec}
retrieves these values, switches {\tt satp} to the kernel page table,
and jumps to {\tt usertrap}, a C function.

The job of {\tt usertrap} is to determine
the cause of the trap, process it, and return
\lineref{kernel/trap.c:/^usertrap/}.
It first changes {\tt stvec} so
that a trap while in the kernel will be handled by
{\tt kernelvec} rather than {\tt uservec}.
It saves the {\tt sepc} register (the saved user program counter)
for future use when returning back to user space.
If the trap is a system call, {\tt usertrap} calls {\tt syscall} to
handle it;
if a device interrupt, {\tt devintr};
if a page fault,  {\tt vmfault};
otherwise it's an exception (e.g., use of an invalid address),
and the kernel kills the faulting process.
The system call path adds four to the saved user program counter
because RISC-V, in the case of a system call,
leaves the program pointer pointing to the {\tt ecall} instruction
but user code needs to resume executing at the subsequent instruction.
{\tt usertrap} checks if the process has been
killed or should yield the CPU (if this trap is a timer interrupt).

The first step in returning to user space is the call to {\tt prepare\_return}
{\lineref{kernel/trap.c:/^prepare/}}.
This function sets up the RISC-V control registers to prepare for a
future trap from user space: setting {\tt stvec}
to {\tt uservec} and preparing the trapframe fields that
{\tt uservec} relies on.
{\tt prepare\_return} sets {\tt sepc} to the previously
saved user program counter.
Finally, {\tt usertrap}
returns back to {\tt userret} in the trampoline page
\lineref{kernel/trampoline.S:/^userret/},
passing back a pointer to the user page table in {\tt a0}.

{\tt userret} switches {\tt satp} to the process's user page table.
Recall that the user page table maps both the trampoline page
and {\tt TRAPFRAME}, but nothing else from the kernel.
The trampoline page mapping at the same
virtual address in user and kernel page tables allows
{\tt userret} to keep executing after changing {\tt satp}.
From this point on, the only data {\tt userret} can use is
the register contents and the content of the trapframe.
{\tt userret} loads the {\tt TRAPFRAME} address into {\tt a0},
restores saved user registers from the trapframe via {\tt a0},
restores the saved user {\tt a0},
and executes {\tt sret} to return to user space.

{\tt uservec} and {\tt userret} are written in assembly
language because it is difficult to write C code to
save or  restore all the registers or survive switching
page tables.

\section{Code: Calling system calls}

User programs call library functions in order to make system calls.
For example, the shell displays a prompt with this function call
(in {\tt user/sh.c}):

\begin{verbatim}
  write(2, "$ ", 2);
\end{verbatim}

Here's the library function, in {\tt user/usys.S}:

\begin{verbatim}
write:
 li a7, SYS_write
 ecall
 ret
\end{verbatim}

The code that the C compiler generates for the function call
loads the three arguments into registers {\tt a0},
{\tt a1}, and {\tt a2}. Then the {\tt write()} function
loads the system call number, {\tt SYS\_write} (16), into {\tt a7}.
The kernel will look at those registers to find out what
system call is intended, and what the arguments are.
The \lstinline{ecall} instruction traps from user space
into the kernel and causes
{\tt uservec},
{\tt usertrap}, and then {\tt syscall} to execute.

At this point, please read \lstinline{kernel/syscall.c} \fileref{kernel/syscall.c},
{\tt sys\_write()} in \lstinline{kernel/sysfile.c} \fileref{kernel/sysfile.c},
and {\tt copyout()}, {\tt copyin()}, and {\tt copyinstr()} in \lstinline{kernel/vm.c} \fileref{kernel/vm.c}.

\indexcode{syscall}
\lineref{kernel/syscall.c:/^syscall/} 
retrieves the system call number from the saved
\texttt{a7} in the trapframe
and uses it to index into {\tt syscalls}
\lineref{kernel/syscall.c:/syscalls/}.
For our example, 
\texttt{a7}
contains 
\indexcode{SYS_write}
\lineref{kernel/syscall.h:/SYS_write/},
resulting in a call to the system call implementation function
\lstinline{sys_write}.

When \lstinline{sys_write} returns,
\lstinline{syscall}
records its return value in
\lstinline{p->trapframe->a0}.
This will cause the original user-space call to 
{\tt write()} to return that value, since the C
calling convention on RISC-V places return values in {\tt a0}.
System calls conventionally return negative numbers to indicate
errors, and zero or positive numbers for success.

\section{Code: System call arguments}

System call arguments start out in the user registers, and
are then moved to the trap frame by the kernel trap code.
The kernel functions
\lstinline{argint},
\lstinline{argaddr},
and
\lstinline{argfd}
retrieve the 
\textit{n} 'th 
system call argument
from the trap frame
as an integer, pointer, or a file descriptor.

Some system calls pass pointers as arguments, and the kernel must use
those pointers to read or write user memory. The {\tt write} system
call, for example, passes the kernel a user-space pointer
to the data to be written.
Such pointers pose
two challenges. First, the user program may be buggy or malicious, and
may pass the kernel an invalid pointer or a pointer intended to trick
the kernel into accessing kernel memory instead of user memory.
Second, the xv6 kernel page table mappings are not the same as the
user page table mappings, so the kernel cannot use ordinary
instructions to load or store from user-supplied addresses.

The kernel implements functions that safely transfer data to and
from user-supplied addresses.
{\tt fetchstr} is an example \lineref{kernel/syscall.c:/^fetchstr/}.
File system calls such as
{\tt exec} use {\tt fetchstr} to retrieve string file-name arguments from user
space.
\lstinline{fetchstr} calls \lstinline{copyinstr}
to do the hard work.

\indexcode{copyinstr}
\lineref{kernel/vm.c:/^copyinstr/} copies up to \lstinline{max} bytes to
\lstinline{dst} from virtual address \lstinline{srcva} in the user page
table \lstinline{pagetable}.
Since \lstinline{pagetable} is {\it not} the current page
table,
\lstinline{copyinstr} uses {\tt walkaddr}
(which calls {\tt walk}) to look up
\lstinline{srcva} in
\lstinline{pagetable}, yielding
physical address \lstinline{pa0}.
The kernel's page table maps all of physical RAM 
at virtual addresses that are equal to the RAM's physical address.
This allows
{\tt copyinstr} to directly copy string bytes from {\tt pa0} to {\tt dst}.
{\tt walkaddr} 
\lineref{kernel/vm.c:/^walkaddr/}
checks that the user-supplied virtual address is part of
the process's user address space, so programs
cannot trick the kernel into reading other memory.
A similar function, {\tt copyout}, copies data from the
kernel to a user-supplied address.

\section{Traps from kernel space}
\label{s:ktraps}

Please read \lstinline{kernel/kernelvec.S} \fileref{kernel/kernelvec.S},
and {\tt kerneltrap()} in \lstinline{kernel/trap.c} \fileref{kernel/trap.c}.

Xv6 handles traps from kernel code in a different way
than traps from user code.
When entering the kernel, {\tt usertrap} points {\tt stvec}
to the assembly code at {\tt kernelvec}
\lineref{kernel/kernelvec.S:/^kernelvec/}.
Since {\tt kernelvec} only executes if
xv6 was already in the kernel, {\tt kernelvec} can rely
on {\tt satp} being set to the kernel page table, and on the
stack pointer referring to a valid kernel stack.
{\tt kernelvec} pushes all 32 registers onto the current stack,
from which it will later restore them
so that the interrupted
kernel code can resume without disturbance.

{\tt kernelvec} saves the registers on the stack of the interrupted
kernel thread, which makes sense because the register values belong to
that thread. This is particularly important if the trap causes a
switch to a different thread -- in that case the trap will actually
return from the stack of the new thread, leaving the interrupted
thread's saved registers safely on its stack.

{\tt kernelvec} jumps to {\tt kerneltrap}
\lineref{kernel/trap.c:/^kerneltrap/} after saving registers.
{\tt kerneltrap} is prepared for just one type of trap:
device interrupts. It calls
{\tt devintr}
\lineref{kernel/trap.c:/^devintr/}
to handle them.
If the trap isn't a device interrupt, it must be an exception,
such as kernel code trying to use an invalid pointer.
This could only be caused by a bug in the kernel code.
The kernel does not have a way to recover in this situation,
so it calls {\tt panic()}, which prints an error message
and then halts.

If {\tt kerneltrap} was called due to a timer interrupt, and a
process's kernel thread is running (as opposed to a scheduler thread),
{\tt kerneltrap} calls {\tt yield} to give other threads a chance to
run. At some point one of those threads will yield, and let our thread
and its {\tt kerneltrap} resume again.
Chapter~\ref{CH:SCHED} explains what happens in {\tt yield}.

When {\tt kerneltrap}'s work is done, it needs to return to whatever
code was interrupted by the trap. Because a {\tt yield} may have
disturbed {\tt sepc} and the previous mode in {\tt sstatus},
{\tt kerneltrap} saves them when it starts. It now restores those
control registers and returns to {\tt kernelvec}
\lineref{kernel/kernelvec.S:/call.kerneltrap$/}.
{\tt kernelvec} pops the saved registers from the stack and
executes {\tt sret}, which copies {\tt sepc} to {\tt pc}
and resumes the interrupted kernel code.

Xv6 sets a CPU's {\tt stvec} to {\tt kernelvec} when that CPU
enters the kernel from user space; you can see this in {\tt usertrap}
\lineref{kernel/trap.c:/DOC: kernelvec/}.
But there's a window of time when the kernel has started executing
but {\tt stvec} is still set to {\tt uservec}, and it's crucial that 
no device interrupt occur during that window.
Luckily the RISC-V always disables interrupts when it starts
to take a trap, and {\tt usertrap} doesn't enable them again until
after it sets {\tt stvec}.

\section{Real world}

The need for trampoline pages could be eliminated if kernel
memory were mapped into every process's user page table (with
\lstinline{PTE_U} clear).
That would
also eliminate the need for a page table switch when trapping from
user space into the kernel. That in turn would allow system call
implementations in the kernel to take advantage of the current
process's user memory being mapped, allowing kernel code to directly
dereference user pointers. Many operating systems have used these ideas to
increase efficiency. Xv6 avoids them in order to reduce the chances of
security bugs in the kernel due to inadvertent use of user pointers,
and to reduce some complexity that would be required to ensure that
user and kernel virtual addresses don't overlap.

\section{Exercises}

\begin{enumerate}

\item Could some or all of the code in {\tt trampoline.S} and
{\tt kernelvec.S} be written in C rather than assembler?

\item Is there a way to eliminate the special {\tt
  TRAPFRAME} page mapping in every user address space? For
  example, could
  {\tt uservec} be modified to simply push the 32 user registers
  onto the kernel stack, or store them in the {\tt proc}
  structure?

\item Could xv6 be modified to eliminate the special {\tt
  TRAMPOLINE} page mapping?

\end{enumerate}

  \end{center}
\caption{Outline of how a trap from user code is handled.}
\label{fig:usertrap}
\end{figure}

Xv6 handles traps differently depending on whether
the trap occurs while
executing in the kernel or in user code. Here is the
story for traps from user code; Section~\ref{s:ktraps}
describes traps from kernel code.

A trap may occur while executing in user space if the
user program makes a
system call ({\tt ecall} instruction), or does something
illegal, or if a device interrupts.
As shown in Figure~\ref{fig:usertrap}, the high-level path of a trap from user space is
{\tt uservec}
\lineref{kernel/trampoline.S:/^uservec/},
then {\tt usertrap}
\lineref{kernel/trap.c:/^usertrap/};
and when the kernel is ready to return,
{\tt usertrap}
returns to
{\tt userret}
\lineref{kernel/trampoline.S:/^userret/}
which executes {\tt sret}
to user space.

% talk about why RISC-V doesn't switch page tables

A major constraint on the design of xv6's trap handling is the fact
that the RISC-V hardware does not switch page tables when it forces a
trap. This means that the trap handler
address in {\tt stvec} must have a valid
mapping in the user page table, since that's the page table in force
when the trap handling code starts executing. Furthermore, xv6's trap
handling code needs to switch to the kernel page table; in order to be
able to continue executing after that switch, the kernel page table
must also have a mapping for the handler pointed to by {\tt stvec}.

Xv6 satisfies these requirements using a \indextext{trampoline} page.
This page contains {\tt uservec}, the xv6 trap handling code
that {\tt stvec} points to. The trampoline page is mapped in
every process's page table at virtual address {\tt 0x3ffffff000}
(called \indexcode{TRAMPOLINE}),
which is the last page in the virtual address space so that it will
be above memory that programs use for themselves.
The trampoline page is mapped at the same virtual address
in the kernel page table. See Figure~\ref{fig:as} and
Figure~\ref{fig:xv6_layout}. Because the trampoline page is mapped in
the user page table, traps can start
executing there in supervisor mode. Because the trampoline page is
mapped at the same address in the kernel address space, the trap handler
can continue to execute after it switches to the kernel page
table.

The code for the {\tt uservec} trap handler is in {\tt trampoline.S}
\lineref{kernel/trampoline.S:/^uservec/}.
When {\tt uservec} starts, all 32 registers contain values owned by
the interrupted user code. These 32 values need to be saved somewhere
in memory, so that later on
the kernel can restore them before returning to user space.
Storing to memory requires use of a register
to hold the store's destination address,
but at this point there are no general-purpose registers available!
Luckily RISC-V provides a helping hand in the
form of the {\tt sscratch} register. The {\tt csrw} instruction at
the start of {\tt uservec} saves {\tt a0} in {\tt
  sscratch}. Now 
{\tt uservec} has
one register ({\tt a0}) to play with.

{\tt uservec}'s next task is to save the 32 user registers.
The kernel allocates, for each process, a page of memory for a
{\tt trapframe} structure that (among other things) has space to
save the 32 user registers
\lineref{kernel/proc.h:/^struct.trapframe/}. Because {\tt satp} still
refers to the user page table, {\tt uservec} needs the trapframe to be
mapped in the user address space. Xv6 maps each process's trapframe
at virtual address {\tt TRAPFRAME} ({\tt 0x3fffffe000}) in that process's user page table;
one page below {\tt TRAMPOLINE}.
Each process's {\tt p->trapframe} 
contains a kernel virtual address for the process's trapframe.

{\tt uservec} sets register {\tt a0} to address {\tt TRAPFRAME}
and saves all the user registers there.
Then it retrieves the user {\tt a0} from {\tt sscratch} and
saves it in the trapframe.

The kernel previously initialized the trapframe to contain some
values useful to {\tt uservec}:
the address of the current process's
kernel stack, the current CPU's hartid, the address of the {\tt usertrap}
function,
and the address of the kernel page table. {\tt uservec}
retrieves these values, switches {\tt satp} to the kernel page table,
and jumps to {\tt usertrap}, a C function.

The job of {\tt usertrap} is to determine
the cause of the trap, process it, and return
\lineref{kernel/trap.c:/^usertrap/}.
It first changes {\tt stvec} so
that a trap while in the kernel will be handled by
{\tt kernelvec} rather than {\tt uservec}.
It saves the {\tt sepc} register (the saved user program counter)
for future use when returning back to user space.
If the trap is a system call, {\tt usertrap} calls {\tt syscall} to
handle it;
if a device interrupt, {\tt devintr};
if a page fault,  {\tt vmfault};
otherwise it's an exception (e.g., use of an invalid address),
and the kernel kills the faulting process.
The system call path adds four to the saved user program counter
because RISC-V, in the case of a system call,
leaves the program pointer pointing to the {\tt ecall} instruction
but user code needs to resume executing at the subsequent instruction.
{\tt usertrap} checks if the process has been
killed or should yield the CPU (if this trap is a timer interrupt).

The first step in returning to user space is the call to {\tt prepare\_return}
{\lineref{kernel/trap.c:/^prepare/}}.
This function sets up the RISC-V control registers to prepare for a
future trap from user space: setting {\tt stvec}
to {\tt uservec} and preparing the trapframe fields that
{\tt uservec} relies on.
{\tt prepare\_return} sets {\tt sepc} to the previously
saved user program counter.
Finally, {\tt usertrap}
returns back to {\tt userret} in the trampoline page
\lineref{kernel/trampoline.S:/^userret/},
passing back a pointer to the user page table in {\tt a0}.

{\tt userret} switches {\tt satp} to the process's user page table.
Recall that the user page table maps both the trampoline page
and {\tt TRAPFRAME}, but nothing else from the kernel.
The trampoline page mapping at the same
virtual address in user and kernel page tables allows
{\tt userret} to keep executing after changing {\tt satp}.
From this point on, the only data {\tt userret} can use is
the register contents and the content of the trapframe.
{\tt userret} loads the {\tt TRAPFRAME} address into {\tt a0},
restores saved user registers from the trapframe via {\tt a0},
restores the saved user {\tt a0},
and executes {\tt sret} to return to user space.

{\tt uservec} and {\tt userret} are written in assembly
language because it is difficult to write C code to
save or  restore all the registers or survive switching
page tables.

\section{Code: Calling system calls}

User programs call library functions in order to make system calls.
For example, the shell displays a prompt with this function call
(in {\tt user/sh.c}):

\begin{verbatim}
  write(2, "$ ", 2);
\end{verbatim}

Here's the library function, in {\tt user/usys.S}:

\begin{verbatim}
write:
 li a7, SYS_write
 ecall
 ret
\end{verbatim}

The code that the C compiler generates for the function call
loads the three arguments into registers {\tt a0},
{\tt a1}, and {\tt a2}. Then the {\tt write()} function
loads the system call number, {\tt SYS\_write} (16), into {\tt a7}.
The kernel will look at those registers to find out what
system call is intended, and what the arguments are.
The \lstinline{ecall} instruction traps from user space
into the kernel and causes
{\tt uservec},
{\tt usertrap}, and then {\tt syscall} to execute.

At this point, please read \lstinline{kernel/syscall.c} \fileref{kernel/syscall.c},
{\tt sys\_write()} in \lstinline{kernel/sysfile.c} \fileref{kernel/sysfile.c},
and {\tt copyout()}, {\tt copyin()}, and {\tt copyinstr()} in \lstinline{kernel/vm.c} \fileref{kernel/vm.c}.

\indexcode{syscall}
\lineref{kernel/syscall.c:/^syscall/} 
retrieves the system call number from the saved
\texttt{a7} in the trapframe
and uses it to index into {\tt syscalls}
\lineref{kernel/syscall.c:/syscalls/}.
For our example, 
\texttt{a7}
contains 
\indexcode{SYS_write}
\lineref{kernel/syscall.h:/SYS_write/},
resulting in a call to the system call implementation function
\lstinline{sys_write}.

When \lstinline{sys_write} returns,
\lstinline{syscall}
records its return value in
\lstinline{p->trapframe->a0}.
This will cause the original user-space call to 
{\tt write()} to return that value, since the C
calling convention on RISC-V places return values in {\tt a0}.
System calls conventionally return negative numbers to indicate
errors, and zero or positive numbers for success.

\section{Code: System call arguments}

System call arguments start out in the user registers, and
are then moved to the trap frame by the kernel trap code.
The kernel functions
\lstinline{argint},
\lstinline{argaddr},
and
\lstinline{argfd}
retrieve the 
\textit{n} 'th 
system call argument
from the trap frame
as an integer, pointer, or a file descriptor.

Some system calls pass pointers as arguments, and the kernel must use
those pointers to read or write user memory. The {\tt write} system
call, for example, passes the kernel a user-space pointer
to the data to be written.
Such pointers pose
two challenges. First, the user program may be buggy or malicious, and
may pass the kernel an invalid pointer or a pointer intended to trick
the kernel into accessing kernel memory instead of user memory.
Second, the xv6 kernel page table mappings are not the same as the
user page table mappings, so the kernel cannot use ordinary
instructions to load or store from user-supplied addresses.

The kernel implements functions that safely transfer data to and
from user-supplied addresses.
{\tt fetchstr} is an example \lineref{kernel/syscall.c:/^fetchstr/}.
File system calls such as
{\tt exec} use {\tt fetchstr} to retrieve string file-name arguments from user
space.
\lstinline{fetchstr} calls \lstinline{copyinstr}
to do the hard work.

\indexcode{copyinstr}
\lineref{kernel/vm.c:/^copyinstr/} copies up to \lstinline{max} bytes to
\lstinline{dst} from virtual address \lstinline{srcva} in the user page
table \lstinline{pagetable}.
Since \lstinline{pagetable} is {\it not} the current page
table,
\lstinline{copyinstr} uses {\tt walkaddr}
(which calls {\tt walk}) to look up
\lstinline{srcva} in
\lstinline{pagetable}, yielding
physical address \lstinline{pa0}.
The kernel's page table maps all of physical RAM 
at virtual addresses that are equal to the RAM's physical address.
This allows
{\tt copyinstr} to directly copy string bytes from {\tt pa0} to {\tt dst}.
{\tt walkaddr} 
\lineref{kernel/vm.c:/^walkaddr/}
checks that the user-supplied virtual address is part of
the process's user address space, so programs
cannot trick the kernel into reading other memory.
A similar function, {\tt copyout}, copies data from the
kernel to a user-supplied address.

\section{Traps from kernel space}
\label{s:ktraps}

Please read \lstinline{kernel/kernelvec.S} \fileref{kernel/kernelvec.S},
and {\tt kerneltrap()} in \lstinline{kernel/trap.c} \fileref{kernel/trap.c}.

Xv6 handles traps from kernel code in a different way
than traps from user code.
When entering the kernel, {\tt usertrap} points {\tt stvec}
to the assembly code at {\tt kernelvec}
\lineref{kernel/kernelvec.S:/^kernelvec/}.
Since {\tt kernelvec} only executes if
xv6 was already in the kernel, {\tt kernelvec} can rely
on {\tt satp} being set to the kernel page table, and on the
stack pointer referring to a valid kernel stack.
{\tt kernelvec} pushes all 32 registers onto the current stack,
from which it will later restore them
so that the interrupted
kernel code can resume without disturbance.

{\tt kernelvec} saves the registers on the stack of the interrupted
kernel thread, which makes sense because the register values belong to
that thread. This is particularly important if the trap causes a
switch to a different thread -- in that case the trap will actually
return from the stack of the new thread, leaving the interrupted
thread's saved registers safely on its stack.

{\tt kernelvec} jumps to {\tt kerneltrap}
\lineref{kernel/trap.c:/^kerneltrap/} after saving registers.
{\tt kerneltrap} is prepared for just one type of trap:
device interrupts. It calls
{\tt devintr}
\lineref{kernel/trap.c:/^devintr/}
to handle them.
If the trap isn't a device interrupt, it must be an exception,
such as kernel code trying to use an invalid pointer.
This could only be caused by a bug in the kernel code.
The kernel does not have a way to recover in this situation,
so it calls {\tt panic()}, which prints an error message
and then halts.

If {\tt kerneltrap} was called due to a timer interrupt, and a
process's kernel thread is running (as opposed to a scheduler thread),
{\tt kerneltrap} calls {\tt yield} to give other threads a chance to
run. At some point one of those threads will yield, and let our thread
and its {\tt kerneltrap} resume again.
Chapter~\ref{CH:SCHED} explains what happens in {\tt yield}.

When {\tt kerneltrap}'s work is done, it needs to return to whatever
code was interrupted by the trap. Because a {\tt yield} may have
disturbed {\tt sepc} and the previous mode in {\tt sstatus},
{\tt kerneltrap} saves them when it starts. It now restores those
control registers and returns to {\tt kernelvec}
\lineref{kernel/kernelvec.S:/call.kerneltrap$/}.
{\tt kernelvec} pops the saved registers from the stack and
executes {\tt sret}, which copies {\tt sepc} to {\tt pc}
and resumes the interrupted kernel code.

Xv6 sets a CPU's {\tt stvec} to {\tt kernelvec} when that CPU
enters the kernel from user space; you can see this in {\tt usertrap}
\lineref{kernel/trap.c:/DOC: kernelvec/}.
But there's a window of time when the kernel has started executing
but {\tt stvec} is still set to {\tt uservec}, and it's crucial that 
no device interrupt occur during that window.
Luckily the RISC-V always disables interrupts when it starts
to take a trap, and {\tt usertrap} doesn't enable them again until
after it sets {\tt stvec}.

\section{Real world}

The need for trampoline pages could be eliminated if kernel
memory were mapped into every process's user page table (with
\lstinline{PTE_U} clear).
That would
also eliminate the need for a page table switch when trapping from
user space into the kernel. That in turn would allow system call
implementations in the kernel to take advantage of the current
process's user memory being mapped, allowing kernel code to directly
dereference user pointers. Many operating systems have used these ideas to
increase efficiency. Xv6 avoids them in order to reduce the chances of
security bugs in the kernel due to inadvertent use of user pointers,
and to reduce some complexity that would be required to ensure that
user and kernel virtual addresses don't overlap.

\section{Exercises}

\begin{enumerate}

\item Could some or all of the code in {\tt trampoline.S} and
{\tt kernelvec.S} be written in C rather than assembler?

\item Is there a way to eliminate the special {\tt
  TRAPFRAME} page mapping in every user address space? For
  example, could
  {\tt uservec} be modified to simply push the 32 user registers
  onto the kernel stack, or store them in the {\tt proc}
  structure?

\item Could xv6 be modified to eliminate the special {\tt
  TRAMPOLINE} page mapping?

\end{enumerate}

  \end{center}
\caption{Outline of how a trap from user code is handled.}
\label{fig:usertrap}
\end{figure}

Xv6 handles traps differently depending on whether
the trap occurs while
executing in the kernel or in user code. Here is the
story for traps from user code; Section~\ref{s:ktraps}
describes traps from kernel code.

A trap may occur while executing in user space if the
user program makes a
system call ({\tt ecall} instruction), or does something
illegal, or if a device interrupts.
As shown in Figure~\ref{fig:usertrap}, the high-level path of a trap from user space is
{\tt uservec}
\lineref{kernel/trampoline.S:/^uservec/},
then {\tt usertrap}
\lineref{kernel/trap.c:/^usertrap/};
and when the kernel is ready to return,
{\tt usertrap}
returns to
{\tt userret}
\lineref{kernel/trampoline.S:/^userret/}
which executes {\tt sret}
to user space.

% talk about why RISC-V doesn't switch page tables

A major constraint on the design of xv6's trap handling is the fact
that the RISC-V hardware does not switch page tables when it forces a
trap. This means that the trap handler
address in {\tt stvec} must have a valid
mapping in the user page table, since that's the page table in force
when the trap handling code starts executing. Furthermore, xv6's trap
handling code needs to switch to the kernel page table; in order to be
able to continue executing after that switch, the kernel page table
must also have a mapping for the handler pointed to by {\tt stvec}.

Xv6 satisfies these requirements using a \indextext{trampoline} page.
This page contains {\tt uservec}, the xv6 trap handling code
that {\tt stvec} points to. The trampoline page is mapped in
every process's page table at virtual address {\tt 0x3ffffff000}
(called \indexcode{TRAMPOLINE}),
which is the last page in the virtual address space so that it will
be above memory that programs use for themselves.
The trampoline page is mapped at the same virtual address
in the kernel page table. See Figure~\ref{fig:as} and
Figure~\ref{fig:xv6_layout}. Because the trampoline page is mapped in
the user page table, traps can start
executing there in supervisor mode. Because the trampoline page is
mapped at the same address in the kernel address space, the trap handler
can continue to execute after it switches to the kernel page
table.

The code for the {\tt uservec} trap handler is in {\tt trampoline.S}
\lineref{kernel/trampoline.S:/^uservec/}.
When {\tt uservec} starts, all 32 registers contain values owned by
the interrupted user code. These 32 values need to be saved somewhere
in memory, so that later on
the kernel can restore them before returning to user space.
Storing to memory requires use of a register
to hold the store's destination address,
but at this point there are no general-purpose registers available!
Luckily RISC-V provides a helping hand in the
form of the {\tt sscratch} register. The {\tt csrw} instruction at
the start of {\tt uservec} saves {\tt a0} in {\tt
  sscratch}. Now 
{\tt uservec} has
one register ({\tt a0}) to play with.

{\tt uservec}'s next task is to save the 32 user registers.
The kernel allocates, for each process, a page of memory for a
{\tt trapframe} structure that (among other things) has space to
save the 32 user registers
\lineref{kernel/proc.h:/^struct.trapframe/}. Because {\tt satp} still
refers to the user page table, {\tt uservec} needs the trapframe to be
mapped in the user address space. Xv6 maps each process's trapframe
at virtual address {\tt TRAPFRAME} ({\tt 0x3fffffe000}) in that process's user page table;
one page below {\tt TRAMPOLINE}.
Each process's {\tt p->trapframe} 
contains a kernel virtual address for the process's trapframe.

{\tt uservec} sets register {\tt a0} to address {\tt TRAPFRAME}
and saves all the user registers there.
Then it retrieves the user {\tt a0} from {\tt sscratch} and
saves it in the trapframe.

The kernel previously initialized the trapframe to contain some
values useful to {\tt uservec}:
the address of the current process's
kernel stack, the current CPU's hartid, the address of the {\tt usertrap}
function,
and the address of the kernel page table. {\tt uservec}
retrieves these values, switches {\tt satp} to the kernel page table,
and jumps to {\tt usertrap}, a C function.

The job of {\tt usertrap} is to determine
the cause of the trap, process it, and return
\lineref{kernel/trap.c:/^usertrap/}.
It first changes {\tt stvec} so
that a trap while in the kernel will be handled by
{\tt kernelvec} rather than {\tt uservec}.
It saves the {\tt sepc} register (the saved user program counter)
for future use when returning back to user space.
If the trap is a system call, {\tt usertrap} calls {\tt syscall} to
handle it;
if a device interrupt, {\tt devintr};
if a page fault,  {\tt vmfault};
otherwise it's an exception (e.g., use of an invalid address),
and the kernel kills the faulting process.
The system call path adds four to the saved user program counter
because RISC-V, in the case of a system call,
leaves the program pointer pointing to the {\tt ecall} instruction
but user code needs to resume executing at the subsequent instruction.
{\tt usertrap} checks if the process has been
killed or should yield the CPU (if this trap is a timer interrupt).

The first step in returning to user space is the call to {\tt prepare\_return}
{\lineref{kernel/trap.c:/^prepare/}}.
This function sets up the RISC-V control registers to prepare for a
future trap from user space: setting {\tt stvec}
to {\tt uservec} and preparing the trapframe fields that
{\tt uservec} relies on.
{\tt prepare\_return} sets {\tt sepc} to the previously
saved user program counter.
Finally, {\tt usertrap}
returns back to {\tt userret} in the trampoline page
\lineref{kernel/trampoline.S:/^userret/},
passing back a pointer to the user page table in {\tt a0}.

{\tt userret} switches {\tt satp} to the process's user page table.
Recall that the user page table maps both the trampoline page
and {\tt TRAPFRAME}, but nothing else from the kernel.
The trampoline page mapping at the same
virtual address in user and kernel page tables allows
{\tt userret} to keep executing after changing {\tt satp}.
From this point on, the only data {\tt userret} can use is
the register contents and the content of the trapframe.
{\tt userret} loads the {\tt TRAPFRAME} address into {\tt a0},
restores saved user registers from the trapframe via {\tt a0},
restores the saved user {\tt a0},
and executes {\tt sret} to return to user space.

{\tt uservec} and {\tt userret} are written in assembly
language because it is difficult to write C code to
save or  restore all the registers or survive switching
page tables.

\section{Code: Calling system calls}

User programs call library functions in order to make system calls.
For example, the shell displays a prompt with this function call
(in {\tt user/sh.c}):

\begin{verbatim}
  write(2, "$ ", 2);
\end{verbatim}

Here's the library function, in {\tt user/usys.S}:

\begin{verbatim}
write:
 li a7, SYS_write
 ecall
 ret
\end{verbatim}

The code that the C compiler generates for the function call
loads the three arguments into registers {\tt a0},
{\tt a1}, and {\tt a2}. Then the {\tt write()} function
loads the system call number, {\tt SYS\_write} (16), into {\tt a7}.
The kernel will look at those registers to find out what
system call is intended, and what the arguments are.
The \lstinline{ecall} instruction traps from user space
into the kernel and causes
{\tt uservec},
{\tt usertrap}, and then {\tt syscall} to execute.

At this point, please read \lstinline{kernel/syscall.c} \fileref{kernel/syscall.c},
{\tt sys\_write()} in \lstinline{kernel/sysfile.c} \fileref{kernel/sysfile.c},
and {\tt copyout()}, {\tt copyin()}, and {\tt copyinstr()} in \lstinline{kernel/vm.c} \fileref{kernel/vm.c}.

\indexcode{syscall}
\lineref{kernel/syscall.c:/^syscall/} 
retrieves the system call number from the saved
\texttt{a7} in the trapframe
and uses it to index into {\tt syscalls}
\lineref{kernel/syscall.c:/syscalls/}.
For our example, 
\texttt{a7}
contains 
\indexcode{SYS_write}
\lineref{kernel/syscall.h:/SYS_write/},
resulting in a call to the system call implementation function
\lstinline{sys_write}.

When \lstinline{sys_write} returns,
\lstinline{syscall}
records its return value in
\lstinline{p->trapframe->a0}.
This will cause the original user-space call to 
{\tt write()} to return that value, since the C
calling convention on RISC-V places return values in {\tt a0}.
System calls conventionally return negative numbers to indicate
errors, and zero or positive numbers for success.

\section{Code: System call arguments}

System call arguments start out in the user registers, and
are then moved to the trap frame by the kernel trap code.
The kernel functions
\lstinline{argint},
\lstinline{argaddr},
and
\lstinline{argfd}
retrieve the 
\textit{n} 'th 
system call argument
from the trap frame
as an integer, pointer, or a file descriptor.

Some system calls pass pointers as arguments, and the kernel must use
those pointers to read or write user memory. The {\tt write} system
call, for example, passes the kernel a user-space pointer
to the data to be written.
Such pointers pose
two challenges. First, the user program may be buggy or malicious, and
may pass the kernel an invalid pointer or a pointer intended to trick
the kernel into accessing kernel memory instead of user memory.
Second, the xv6 kernel page table mappings are not the same as the
user page table mappings, so the kernel cannot use ordinary
instructions to load or store from user-supplied addresses.

The kernel implements functions that safely transfer data to and
from user-supplied addresses.
{\tt fetchstr} is an example \lineref{kernel/syscall.c:/^fetchstr/}.
File system calls such as
{\tt exec} use {\tt fetchstr} to retrieve string file-name arguments from user
space.
\lstinline{fetchstr} calls \lstinline{copyinstr}
to do the hard work.

\indexcode{copyinstr}
\lineref{kernel/vm.c:/^copyinstr/} copies up to \lstinline{max} bytes to
\lstinline{dst} from virtual address \lstinline{srcva} in the user page
table \lstinline{pagetable}.
Since \lstinline{pagetable} is {\it not} the current page
table,
\lstinline{copyinstr} uses {\tt walkaddr}
(which calls {\tt walk}) to look up
\lstinline{srcva} in
\lstinline{pagetable}, yielding
physical address \lstinline{pa0}.
The kernel's page table maps all of physical RAM 
at virtual addresses that are equal to the RAM's physical address.
This allows
{\tt copyinstr} to directly copy string bytes from {\tt pa0} to {\tt dst}.
{\tt walkaddr} 
\lineref{kernel/vm.c:/^walkaddr/}
checks that the user-supplied virtual address is part of
the process's user address space, so programs
cannot trick the kernel into reading other memory.
A similar function, {\tt copyout}, copies data from the
kernel to a user-supplied address.

\section{Traps from kernel space}
\label{s:ktraps}

Please read \lstinline{kernel/kernelvec.S} \fileref{kernel/kernelvec.S},
and {\tt kerneltrap()} in \lstinline{kernel/trap.c} \fileref{kernel/trap.c}.

Xv6 handles traps from kernel code in a different way
than traps from user code.
When entering the kernel, {\tt usertrap} points {\tt stvec}
to the assembly code at {\tt kernelvec}
\lineref{kernel/kernelvec.S:/^kernelvec/}.
Since {\tt kernelvec} only executes if
xv6 was already in the kernel, {\tt kernelvec} can rely
on {\tt satp} being set to the kernel page table, and on the
stack pointer referring to a valid kernel stack.
{\tt kernelvec} pushes all 32 registers onto the current stack,
from which it will later restore them
so that the interrupted
kernel code can resume without disturbance.

{\tt kernelvec} saves the registers on the stack of the interrupted
kernel thread, which makes sense because the register values belong to
that thread. This is particularly important if the trap causes a
switch to a different thread -- in that case the trap will actually
return from the stack of the new thread, leaving the interrupted
thread's saved registers safely on its stack.

{\tt kernelvec} jumps to {\tt kerneltrap}
\lineref{kernel/trap.c:/^kerneltrap/} after saving registers.
{\tt kerneltrap} is prepared for just one type of trap:
device interrupts. It calls
{\tt devintr}
\lineref{kernel/trap.c:/^devintr/}
to handle them.
If the trap isn't a device interrupt, it must be an exception,
such as kernel code trying to use an invalid pointer.
This could only be caused by a bug in the kernel code.
The kernel does not have a way to recover in this situation,
so it calls {\tt panic()}, which prints an error message
and then halts.

If {\tt kerneltrap} was called due to a timer interrupt, and a
process's kernel thread is running (as opposed to a scheduler thread),
{\tt kerneltrap} calls {\tt yield} to give other threads a chance to
run. At some point one of those threads will yield, and let our thread
and its {\tt kerneltrap} resume again.
Chapter~\ref{CH:SCHED} explains what happens in {\tt yield}.

When {\tt kerneltrap}'s work is done, it needs to return to whatever
code was interrupted by the trap. Because a {\tt yield} may have
disturbed {\tt sepc} and the previous mode in {\tt sstatus},
{\tt kerneltrap} saves them when it starts. It now restores those
control registers and returns to {\tt kernelvec}
\lineref{kernel/kernelvec.S:/call.kerneltrap$/}.
{\tt kernelvec} pops the saved registers from the stack and
executes {\tt sret}, which copies {\tt sepc} to {\tt pc}
and resumes the interrupted kernel code.

Xv6 sets a CPU's {\tt stvec} to {\tt kernelvec} when that CPU
enters the kernel from user space; you can see this in {\tt usertrap}
\lineref{kernel/trap.c:/DOC: kernelvec/}.
But there's a window of time when the kernel has started executing
but {\tt stvec} is still set to {\tt uservec}, and it's crucial that 
no device interrupt occur during that window.
Luckily the RISC-V always disables interrupts when it starts
to take a trap, and {\tt usertrap} doesn't enable them again until
after it sets {\tt stvec}.

\section{Real world}

The need for trampoline pages could be eliminated if kernel
memory were mapped into every process's user page table (with
\lstinline{PTE_U} clear).
That would
also eliminate the need for a page table switch when trapping from
user space into the kernel. That in turn would allow system call
implementations in the kernel to take advantage of the current
process's user memory being mapped, allowing kernel code to directly
dereference user pointers. Many operating systems have used these ideas to
increase efficiency. Xv6 avoids them in order to reduce the chances of
security bugs in the kernel due to inadvertent use of user pointers,
and to reduce some complexity that would be required to ensure that
user and kernel virtual addresses don't overlap.

\section{Exercises}

\begin{enumerate}

\item Could some or all of the code in {\tt trampoline.S} and
{\tt kernelvec.S} be written in C rather than assembler?

\item Is there a way to eliminate the special {\tt
  TRAPFRAME} page mapping in every user address space? For
  example, could
  {\tt uservec} be modified to simply push the 32 user registers
  onto the kernel stack, or store them in the {\tt proc}
  structure?

\item Could xv6 be modified to eliminate the special {\tt
  TRAMPOLINE} page mapping?

\end{enumerate}
